\section{Admissible physical curve growth and the clearance consequence}
\label{sec:uc-growth}
We now separate the mathematical application of a published combinatorial
lemma from the inherited global spectral and roof-pressure inputs. The
application below is a theorem of this revision, not a claim that a
successful build has certified the billiard dynamics.

\begin{theorem}[The required angular-source growth input is satisfied]
\label{thm:uc-growth}
For the original triangular circular Lorentz family, choose the per-radius
unstable graph class
\[
 \mathcal W_R^u(B_{\rm cv})=\{C:\ C\text{ is a homogeneous }C^2\text{ graph},\
           \kappa_R\le d\varphi/ds\le U,\ |K_C|\le B_{\rm cv}\},
\]
where $B_{\rm cv}$ satisfies \eqref{eq:uc-B}. There exist common $k_0,\ell_*>0$
and $C_g<\infty$ such that, for every $R$, every
$C\in\mathcal W_R^u(B_{\rm cv})$ with $|C|\le\ell_*$ and every integer $l\ge0$,
\begin{equation}\label{eq:uc-growth-conclusion}
 \sum_{V\in\mathcal G_l^+(C)}\sup_V J_V(T_R^{-l})\le C_g.
\end{equation}
The partition is obtained by the actual regular homogeneous image branches
and a fixed short-length subdivision. It covers arclength-almost all of
$C$ before pushforward. It contains no added reflection at a clearance mark.
All hypotheses of Input~\ref{input:ag-growth}, for this chosen class, are
therefore satisfied.
\end{theorem}
\begin{proof}
Fix $k_0$ by \eqref{eq:uc-k0}; choose $\ell_*$ no larger than the
length from Lemma~\ref{lem:uc-central}, the short-curve length needed
for its preliminary relative-angle argument, and the length in
\eqref{eq:ag-length-choice}. All choices are made once for the
radius family, before the collision count or any outer limit.

Theorem~\ref{thm:uc-curvature} verifies invariance of the homogeneous
$C^2$ curve class in the literal geometric-curvature sense of (H2) in
\cite{DZUniform}. It contains the original first-image source graphs;
this membership is now proved, not inferred from their being $C^2$.
Equation~\eqref{eq:uc-adapted} gives the common adapted hyperbolicity
and a uniform endpoint metric comparison. Theorem~\ref{thm:uc-expansion}
verifies (H3.A) and (H3.B), with the explicit permissible threshold
$\zeta_0=1/2$. Its estimate uses suprema of contraction, so it is at
least as strong as the one-step hypothesis used there. Theorem~
\ref{thm:uc-distortion} gives the arclength and invariant-measure parts
of (H4) on homogeneous branches. The logarithmic and ratio forms are
converted on the one fixed short length. None of these verifications
uses the physical-word pressure input.

For completeness, singular-null coverage is checked on these curves.
At each forward step there are only finitely many physical tangency
cuts on a short graph, by the positive-foot monotonicity above. Its
regular image has strictly monotone $\varphi_1$, so homogeneous
boundaries have at most one preimage per first-hit branch per strip.
The additional length cuts are countable. Starting with a countable
collection, iteration thus excludes a countable collection of points
for each finite $l$. Each regular restriction is injective. No positive
arclength part of a graph lies in a tangency or homogeneous-boundary
preimage. This proves the asserted null coverage without treating a
bulk-null set as a zero endpoint trace in any pairing construction.

We use precisely \cite[Lemma 3.2(b), p.~14]{DZUniform}: for the stable
class and the homogeneous inverse partition $\mathcal G_l(W)$, with
long pieces subdivided at the fixed short scale, the sum of sup
arclength contraction Jacobians is at most a constant depending on
the common curve, expansion and distortion constants, uniformly in
$W$, $T$ and $l$. The published proof specifies (H2)--(H4) as its
requirements; it is not a theorem about all cylinder suprema or roof
Fourier transforms. Apply this statement to
$\mathcal W_R^s(B_{\rm cv})=\mathcal I\mathcal W_R^u(B_{\rm cv})$ and reverse it by
$\mathcal I T_R\mathcal I=T_R^{-1}$. Since $\mathcal I$ is an
arclength isometry, its Jacobian conversion is exact and gives
\eqref{eq:uc-growth-conclusion}. At $l=0$ the sum is one.
One may apply the lemma separately to each $R$: its constant depends
only on the common estimates just proved. A common Banach-space
identification across radii, or a perturbation theorem for variable
boundary lengths, is not required for this curve-sum conclusion.
Further fixed shortening of the scale is included in $C_g$.
\end{proof}

\begin{corollary}[Complete physical two-clearance bound with the input discharged]
\label{cor:uc-clearance}
Retain the genuine one-copy cross-flight coarea and augmented-itinerary
uniqueness theorems of Section~\ref{sec:xf-jacobi}, with their stated
physical scopes. In the original Lorentz family the bounds
\eqref{eq:ag-kernel}, \eqref{eq:ag-roots}, \eqref{eq:ag-pair},
\eqref{eq:ag-factorial} and \eqref{eq:ag-exact-source} now follow
without an unverified homogeneous curve-growth premise and without
physical-word pressure. In particular, for every fixed
$0<u_1\le\delta_*$ and every $h\ge1$,
\begin{equation}\label{eq:uc-kernel}
 \sum_{\sigma=\pm1}\int_0^{2\pi}
    N_I(T_R^h X_R(\vartheta,\sigma,u_1))\,d\vartheta
                 \le A_*C_g C_t\,|I|.
\end{equation}
The first-image source is taken after exactly one actual collision, and
there are $h-1$ further collisions before applying the second mark.
All exact-record and original-roof Borel selectors between zero and one
may be kept inside this positive inequality. The full one-copy root sum
is bounded in clearance essential supremum uniformly in $R$ and $h$.
\end{corollary}
\begin{proof}
Theorem~\ref{thm:uc-growth} supplies exactly the class, norm, partition,
null-set and uniformity requirements of Input~\ref{input:ag-growth}.
Insert it in the preserved proofs of Theorems~\ref{thm:ag-kernel} and
\ref{thm:ag-roots} and Corollary~\ref{cor:ag-record}. The angular source
has mass $4\pi$; the collision flux factor is $1/(4\pi R)$, and
conditioning on the original section contributes exactly $1/c_*$.
No further angular sign factor is inserted into $A_*C_gC_t$.
The factorial conclusions and their constants are inherited consequences;
the new assertion is the discharge of their curve-growth dependency.
A null exceptional set can be chosen for all integer $h$ at a fixed
radius. No null set common to uncountably many radii is claimed.
\end{proof}

\subsection*{The distinction from the original roofwise theorem}
Theorem~\ref{thm:uc-growth} is a complete curve-sum conclusion and
Corollary~\ref{cor:uc-clearance} sums all regular physical two-mark
branches, including arbitrarily large separation and arbitrarily small
positive incidences. They remove a genuine hypothesis from that part of
the proof chain. They do not turn a clearance density into a roof density.
The count mean $S/R$ is still first order, and the original-clock
pushforward can have folds, critical and one-sided boundary contributions.
In particular no factor $m^{-2}$ occurs in \eqref{eq:uc-kernel}.

The original target remains, for every fixed central parameter $M$,
\begin{equation}\label{eq:uc-open-limit}
 \lim_{B\to\infty}\limsup_{m\to\infty}\sup_{R,n,k}m^2
 \esssup_{t\in\mathcal C_{m,R,n,k}(M)}
 \bigl(\mathsf U^{\rm low}_{B,m}(t)
                   -K_B P^{\rm ng}_{r_B,m}(t)\bigr)_+=0.
\end{equation}
Here $B$ is the inherited outer parameter. The curvature allowance $B_{\rm cv}$
is fixed permanently by \eqref{eq:uc-B}. All original
outer widths are fixed before the collision limsup. Their prescribed
$B\to\infty$ limit is not used to tune the curve partition.
No result of this revision establishes \eqref{eq:uc-open-limit}, the full
raw pointwise LLT, or its final four-coordinate inversion. The remaining
minimal obligations along this route are the normalized first-order
physical-roof source and complete chart/boundary sum, the incidence-only
low-load source or valid comparison capacity, and the stated original
spectral, arithmetic, pairing and normalization interfaces. The new
one-flight complexity and curve-sum proof does not certify any stronger
whole-cylinder or roofwise pressure statement still used elsewhere.

